\section{Related literature}
\label{sec:literature}

\paragraph{The theory-based view.}
The foundational claim is that theories do more than summarize evidence: they
direct attention toward causal possibilities, support counterfactual
reasoning, and guide experimentation \citep{FelinZenger2017}.  The
\emph{Strategy Science} collection on theory-based decisions substantially
developed the view \citep{FelinGambardellaZenger2024}:
\citet{CamuffoEtAl2024} formalize the choice of which theory to test and show
that the value of experimentation grows with the number and uncertainty of
candidate theories; \citet{ChavdaGansStern2024} model search guided by a
given causal theory; \citet{EhrigZenger2024} translate theories into
awareness structures and show how partial disclosure secures resources and
rents; \citet{Sorenson2024} and \citet{AdnerLevinthal2024} probe the
mechanisms and the limits of theory-guided experimentation; and
\citet{OttHannah2024} document empirically how entrepreneurs develop a
poorly formed thesis into an articulated causal model.  A parallel empirical
program shows that training entrepreneurs to formulate and test theories
improves performance \citep{CamuffoEtAl2020}.  Related conceptual work
presses the same open question this paper addresses: \citet{ZellwegerZenger2023}
treat entrepreneurs as scientists generating beliefs under uncertainty, and
\citet{FelinHolweg2024} argue that forward-looking causal theorizing is
precisely what data-driven prediction cannot supply.  Throughout, the
formal treatments condition on the theory, the attribute list, or the
awareness advantage being present.  This paper supplies the front end those
models presuppose, and its equilibrium analysis shows when the supplied front
end matters: heterogeneous causal languages, and the performance differences
they support, persist exactly when the strategic and environmental conditions
identified below fail to erase them.

\paragraph{Unawareness and growing awareness.}
Standard state-space models cannot represent nontrivial unawareness
\citep{DekelLipmanRustichini1998}; generalized structures represent multiple
awareness levels and interactive beliefs
\citep{HeifetzMeierSchipper2006,HeifetzMeierSchipper2013}.  On belief
revision under expansion, reverse Bayesianism preserves relative odds as the
state space grows \citep{KarniViero2013}, with extensions to awareness of
unawareness \citep{KarniViero2017}, intertemporal settings \citep{Viero2021},
unforeseen evidence \citep{Piermont2021}, and learning under incomplete
awareness \citep{GrantMeneghelTourky2022}.  \citet{Schipper2021}
distinguishes learning, which discards conceived possibilities, from
discovery, which adds unconceived ones, and shows how play can move players
to richer games; \citet{Schipper2025} shows that beliefs about novelty can be
held without representing the novel objects---a precedent for the delabeled
constructs used here.  Relative to this literature the paper endogenizes what
those models take as given: the expansion event itself.  Experience generates
the defect, the defect localizes the question, a fallible technology produces
the new distinction, and the bridge prior problem---which reverse Bayesianism
constrains but does not solve---is embedded in an explicit account of what
the generating evidence can and cannot warrant.  In strategy,
\citet{BryanRyallSchipper2022} bring unawareness into value capture and show
that awareness itself earns rents; the present paper models where the
awareness advantage originates, which that paper leaves open.

\paragraph{Causal learning with a fixed vocabulary.}
Causal graphical models separate intervention from association
\citep{Pearl2009,SpirtesEtAl2000}; \citet{Ryall2009} brings this machinery to
strategy, modeling imitation as Bayesian inference over operating structures
in which passive observation narrows candidates to an
observational-equivalence class and experimentation may discriminate among
survivors.  Active-learning methods select interventions to split a
represented equivalence class \citep{HeGeng2008}.  These models are the
benchmark for the paper's second learning tier: inference over structures
expressible in a fixed vocabulary.  The paper begins where they stop---when
no expressible structure is adequate---and its nonidentification results
extend the equivalence-class logic to the choice among vocabulary repairs:
observationally equivalent worlds can demand different distinctions, so no
truth-blind repair rule is uniformly correct.

\paragraph{Subjective rationality and self-confirming equilibrium.}
The back end follows \citet{KalaiLehrer1993b} and \citet{FudenbergLevine1993}:
beliefs correct on the path of play need not be correct off it, and
subjectively optimal behavior against such beliefs can persist.
\citet{Ryall2003} showed that competitive advantage can rest entirely on
rivals' self-confirming false theories.  Those models fix the hypothesis
space; convergence results in the rational-learning tradition require a grain
of truth---the truth absolutely continuous with respect to beliefs
\citep{KalaiLehrer1993a}.  Under unawareness the grain-of-truth condition can
fail at the level of language: the truth may be inexpressible.  This paper
supplies the missing dynamics.  Statistical dissonance is the detector of
language-level failure, awareness refinement restores the grain of truth on
the path of play, and the fixed-language models are recovered exactly as the
special case in which awareness never changes.  The paper's absorbing
configurations extend self-confirming equilibrium with two epistemic
conditions---the language survives its own evidence, and no question is worth
its price---and it is these conditions, not the strategic ones, that carry
the new results.

\paragraph{Statement of contribution.}
Existing models show how theories guide learning, how Bayesian agents
discriminate among represented causal structures, and how play can reveal
previously unavailable actions.  This paper models a different transition.
An agent's causal language pools objectively different situations;
strategic experience endogenously reveals an action-relevant failure of that
pooling and thereby creates a localized open question; costly inquiry
activates a fallible technology whose output is not measurable in the agent's
prior awareness; only after that output is formulated does the agent
construct a prior and undertake Bayesian judgment, under an explicit account
of how much confidence the generating experience licenses.  Because actions
depend on theories, play can either create the contrasts that generate
questions or suppress them---so the model also delivers the population-level
answer: when theory heterogeneity and performance differences persist, and
when the market erases them.
